\documentclass{article}
\usepackage[utf8]{inputenc}
\usepackage{amsmath,amssymb}
\usepackage[most]{tcolorbox}
\usepackage{geometry}
\geometry{margin=2.5cm}

\newcommand{\step}[1]{\par\noindent\hangindent=3.5em\hangafter=1%
  \makebox[3.5em][l]{\textbf{#1.}}}
\newcommand{\steptag}[1]{\hfill{\small\textsf{[#1]}}}

\newtcolorbox{proofbox}[1]{
  colback=gray!5, colframe=gray!60,
  fonttitle=\bfseries, title={#1},
  breakable, enhanced,
  left=4pt, right=4pt, top=4pt, bottom=4pt,
  before skip=10pt, after skip=10pt
}

\begin{document}

\begin{proofbox}{For all reals K\_R, L, K\_C >= 0 there is a universal delta...}
\step{1} For all reals K\_R, L, K\_C >= 0 there is a universal delta\_R(K\_R,L,K\_C) in (0,2\textasciicircum{}(-16)] such that no finite exact signed idempotent P with 0 < delta(P) <= delta\_R admits full row-point fibers represented by v and f, a real A >= 4 and a point q of the row polytope K(P) with sqrt(delta(P))/2 <= ||q - p\_v||\_1 <= 2*sqrt(delta(P)) and ||p\_f - p\_v + A*(q - p\_v)||\_1 <= K\_R*delta(P), and an affine chi with chi(p\_v) = 0, chi(q) = 1, and |chi(x) - chi(y)| <= ||x - y||\_1/||q - p\_v||\_1 for all x,y in K(P), such that sum\_\{Q: |chi(p\_Q)| > L\} max(c\_Q, 0) <= K\_C*delta(P), where c\_Q = sum\_\{j in Q\} P\_\{vj\}. \steptag{validated} \\[6pt]
\step{1.1} Fix K\_R,L,K\_C >= 0 and set B := K\_C+1+(K\_C+K\_R+1)/4, H := 2L+6B, and delta\_R := min\{2\textasciicircum{}(-16), 1/(4H\textasciicircum{}2)\}. Then B >= 5/4 and H > 0, so delta\_R depends only on the three fixed real parameters, is independent of the dimension and of P, and belongs to (0,2\textasciicircum{}(-16)]. \steptag{validated} \\[4pt]
\step{1.2} For these constants, every putative configuration satisfying all hypotheses of node 1 with delta := delta(P) in (0,delta\_R], tau := sqrt(delta), D := q-p\_v, and s := ||D||\_1 necessarily satisfies 1 <= tau*[2L+2(2+4delta)B]. \steptag{validated} \\[6pt]
\step{1.2.1} Let I be the finite index set and write q as a convex combination q = sum\_i lambda\_i p\_i. By def-signed-idempotent every row p\_i has total mass one, and by def-negative-mass its negative mass nu(p\_i) := sum\_j max(-P\_ij,0) is at most delta. Coordinatewise convexity of x -> max(-x,0) gives q*1 = 1 and nu(q) <= sum\_i lambda\_i nu(p\_i) <= delta. Consequently, for every coordinate subset U, q(U) >= -delta and p\_f(U) >= -delta. \steptag{validated} \\[4pt]
\step{1.2.2} The scale window gives s >= tau/2 > 0, hence q != p\_v. For every full row-point fiber Q define d\_Q := sum\_\{j in Q\}(q\_j-P\_vj). Applying the registered external lem-hx-transverse-moment-identity with q0=p\_v, q1=q, and the given affine chi gives the exact unit moment sum\_Q d\_Q*chi(p\_Q) = 1. \steptag{validated} \\[4pt]
\step{1.2.3} Every actual row p\_i has ||p\_i||\_1 = 1+2nu(p\_i) <= 1+2delta, so any two row points have l1-distance at most 2+4delta. Thus the Lipschitz hypothesis and chi(p\_v)=0 give |chi(p\_Q)| <= (2+4delta)/s for every fiber Q. For the core C := \{Q: |chi(p\_Q)| <= L\}, grouping coordinates cannot increase variation, so sum\_\{Q in C\}|d\_Q| <= sum\_Q|d\_Q| <= ||q-p\_v||\_1 = s, and therefore the absolute core contribution is at most Ls. \steptag{validated} \\[4pt]
\step{1.2.4} Let T := \{Q: |chi(p\_Q)| > L\}, split it as T\_- := \{Q in T: d\_Q<0\}, T\_+ := \{Q in T: d\_Q>0\}, and T\_0 := \{Q in T: d\_Q=0\}, and put c\_Q := sum\_\{j in Q\}P\_vj. The tail cap, the negative-mass subset budgets, A>=4, and R := p\_f-p\_v+A(q-p\_v) with ||R||\_1 <= K\_R*delta imply sum\_\{Q in T\}|d\_Q| <= [K\_C+1+(K\_C+K\_R+1)/4]*delta = B*delta. \steptag{validated} \\[6pt]
\step{1.2.4.1} Let U\_- be the union of the coordinate fibers in T\_-. Since d\_Q = sum\_\{j in Q\}(q\_j-P\_vj) = q(Q)-c\_Q is negative on T\_-, one has sum\_\{Q in T\_-\}|d\_Q| = c(U\_-)-q(U\_-). Now c(U\_-) <= sum\_\{Q in T\_-\}max(c\_Q,0) <= K\_C*delta by the tail cap, while q(U\_-) >= -nu(q) >= -delta. Hence sum\_\{Q in T\_-\}|d\_Q| <= (K\_C+1)*delta. \steptag{validated} \\[4pt]
\step{1.2.4.2} Let U\_+ be the union of the coordinate fibers in T\_+ and D\_+ := sum\_\{Q in T\_+\}d\_Q = sum\_\{Q in T\_+\}|d\_Q|. Summing R := p\_f-p\_v+A(q-p\_v) over U\_+ gives R(U\_+) = p\_f(U\_+)-c(U\_+)+A*D\_+, hence A*D\_+ = c(U\_+)+R(U\_+)-p\_f(U\_+). The tail cap gives c(U\_+) <= K\_C*delta, the residual bound gives R(U\_+) <= ||R||\_1 <= K\_R*delta, and def-negative-mass gives -p\_f(U\_+) <= nu(p\_f) <= delta. Thus A*D\_+ <= (K\_C+K\_R+1)*delta and, because A>=4, sum\_\{Q in T\_+\}|d\_Q| <= (K\_C+K\_R+1)*delta/4. \steptag{validated} \\[4pt]
\step{1.2.4.3} If sum\_\{Q in T\_-\}|d\_Q| <= (K\_C+1)*delta and sum\_\{Q in T\_+\}|d\_Q| <= (K\_C+K\_R+1)*delta/4, then, because T\_-, T\_+, T\_0 are disjoint and every Q in T\_0 has d\_Q=0, one has sum\_\{Q in T\}|d\_Q| <= (K\_C+1)*delta+(K\_C+K\_R+1)*delta/4 = B*delta. \steptag{validated} \\[4pt]
\step{1.2.5} Whenever the same putative configuration has the unit moment sum\_Q d\_Q*chi(p\_Q)=1, the core bounds |chi(p\_Q)|<=L and sum\_C|d\_Q|<=s, and the tail bounds |chi(p\_Q)|<=(2+4delta)/s and sum\_T|d\_Q|<=B*delta, the triangle inequality yields 1 <= Ls+(2+4delta)B*delta/s. The actor window s<=2tau and s>=tau/2, together with delta=tau\textasciicircum{}2, then gives 1 <= tau*[2L+2(2+4delta)B]. \steptag{validated} \\[4pt]
\step{1.3} For the constants above and every 0 < delta <= delta\_R, the numerical inequality 1 <= sqrt(delta)*[2L+2(2+4delta)B] is impossible: delta <= 2\textasciicircum{}(-16) gives 2+4delta < 3, while delta <= 1/(4H\textasciicircum{}2) gives sqrt(delta)*H <= 1/2; since B>0, the displayed right side is strictly less than sqrt(delta)*(2L+6B) = sqrt(delta)*H <= 1/2. \steptag{validated} \\[4pt]
\end{proofbox}

\end{document}
